\documentclass[11pt]{article}
\usepackage[utf8]{inputenc}
\usepackage[T1]{fontenc}
\usepackage{amsmath, amssymb, amsthm}
\usepackage{geometry}
\geometry{a4paper, margin=1in}

\usepackage{xcolor}

\begin{document}

\title{\textbf{Note on Some Pseudo-Elliptic Integrals} \\ \large{\textnormal{by E. Goursat}}}
\author{}
\date{Session of March 16, 1887 \\ \textit{\small{Originally published in the Bulletin de la S. M. F., tome 15 (1887), p. 106-120}}}
\maketitle

\theoremstyle{definition}
\newtheorem{definition}{Definition}
\newtheorem*{theorem*}{Theorem}
\newtheorem{theorem}{Theorem}
\newcommand{\Fcal}{\mathfrak{F}}

\section*{1}

Consider a differential of the form
\begin{equation}
f(x)\frac{dx}{\sqrt{Ax^{2}+Bx+C}},
\end{equation}
where $f(x)$ denotes an arbitrary rational function and $Ax^{2}+Bx+C$ is a second-degree trinomial that is not a perfect square, which may, as a special case, reduce to a linear expression. Let us pose the substitution
\begin{equation}
x=\frac{at^{2}+bt+c}{a^{\prime}t^{2}+b^{\prime}t+c^{\prime}}=\varphi(t);
\end{equation}
the numerator and the denominator of $\varphi(t)$ are assumed to be coprime, and at least one is of the second degree. By this substitution, the differential (1) becomes
\begin{equation}
\Fcal(t)\frac{dt}{\sqrt{R(t)}}
\end{equation}
where
\begin{align*}
\Fcal(t) &= f\left(\frac{at^{2}+bt+c}{a^{\prime}t^{2}+b^{\prime}t+c^{\prime}}\right)\frac{Lt^{2}-2Nt+M}{(a^{\prime}t^{2}+b^{\prime}t+c^{\prime})^{2}} \\
L &= ab^{\prime}-ba^{\prime}, \quad M=bc^{\prime}-cb^{\prime}, \quad N=ca^{\prime}-ac^{\prime}, \\
R(t) &= A(at^{2}+bt+c)^{2}+B(at^{2}+bt+c)(a^{\prime}t^{2}+b^{\prime}t+c^{\prime})+C(a^{\prime}t^{2}+b^{\prime}t+c^{\prime})^{2}.
\end{align*}
The two roots $t_{1}$, $t_{2}$ of the equation $\varphi(t)=\lambda$ satisfy, for any $\lambda$, the involution relation
\begin{equation*}
L t_{1}t_{2}-N(t_{1}+t_{2})+M=0,
\end{equation*}
obtained by eliminating $\lambda$ between the two equations
\begin{align*}
a(t_{1}+t_{2})+b-\lambda[a^{\prime}(t_{1}+t_{2})+b^{\prime}] &= 0, \\
at_{1}t_{2}-c-\lambda[a^{\prime}t_{1}t_{2}-c^{\prime}] &= 0.
\end{align*}
It follows that if, in $\varphi(t)$, we make the linear substitution
\begin{equation}
Ltz - N(t+z) + M = 0, \quad \text{or} \quad t = \frac{Nz-M}{Lz-N},
\end{equation}
we will have identically $\varphi(t)=\varphi(z)$. Indeed, one can verify without any difficulty the identities below:
\begin{align*}
at^{2}+bt+c &= (N^{2}-LM)\frac{az^{2}+bz+c}{(Lz-N)^{2}}; \\
a^{\prime}t^{2}+b^{\prime}t+c^{\prime} &= (N^{2}-LM)\frac{a^{\prime}z^{2}+b^{\prime}z+c^{\prime}}{(Lz-N)^{2}}; \\
Lt^{2}-2Nt+M &= (LM-N^{2})\frac{Lz^{2}-2Nz+M}{(Lz-N)^{2}}.
\end{align*}
We therefore have
\begin{align*}
\frac{dt}{\sqrt{R(t)}} &= \frac{(LM-N^{2})dz}{\sqrt{(N^{2}-LM)^{2}R(z)}} = \frac{dz}{\sqrt{R(z)}}, \\
f\left(\frac{at^{2}+bt+c}{a^{\prime}t^{2}+b^{\prime}t+c^{\prime}}\right) &= f\left(\frac{az^{2}+bz+c}{a^{\prime}z^{2}+b^{\prime}z+c^{\prime}}\right); \\
\frac{Lt^{2}-2Nt+M}{(a^{\prime}t^{2}+b^{\prime}t+c^{\prime})^2} &= -\frac{Lz^{2}-2Nz+M}{(a^{\prime}z^{2}+b^{\prime}z+c^{\prime})^2}.
\end{align*}
It follows that
\begin{equation*}
\Fcal(t) = f\left(\frac{at^{2}+bt+c}{a^{\prime}t^{2}+b^{\prime}t+c^{\prime}}\right)\frac{Lt^{2}-2Nt+M}{(a^{\prime}t^{2}+b^{\prime}t+c^{\prime})^{2}} = -f\left(\frac{az^{2}+bz+c}{a^{\prime}z^{2}+b^{\prime}z+c^{\prime}}\right)\frac{Lz^{2}-2Nz+M}{(a^{\prime}z^{2}+b^{\prime}z+c^{\prime})^{2}} = -\Fcal(z).
\end{equation*}
The differential (3) therefore enjoys the following property: \\

\textit{There exists an involutive linear substitution $t=\frac{Nz-M}{Lz-N}$, which permutes the roots of the equation $R(t)=0$ in pairs, such that we identically have $\Fcal\left(\frac{Nz-M}{Lz-N}\right)=-\Fcal(z)$.}

\section*{2}

Conversely, let $R(t)$ be a fourth-degree polynomial having four distinct roots $a, b, c, d$, and let
\begin{equation*}
Lt_{1}t_{2}-N(t_{1}+t_{2})+M=0
\end{equation*}
be an involutive linear substitution which permutes the roots $a, b, c, d$ in pairs. If a rational function $\Fcal(t)$ is such that we identically have
\begin{equation*}
\Fcal(t_{1})+\Fcal(t_{2})=0,
\end{equation*}
the integral
\begin{equation*}
\int\Fcal(t)\frac{dt}{\sqrt{R(t)}}
\end{equation*}
is a pseudo-elliptic integral. Let us indeed denote by $\alpha$ and $\beta$ the double points of the preceding involution; the relation between $t_1$ and $t_2$ can be written
\begin{equation*}
\frac{t_{1}-\alpha}{t_{1}-\beta}+\frac{t_{2}-\alpha}{t_{2}-\beta}=0.
\end{equation*}
Let's make the change of variable,
\begin{equation*}
\frac{t-\alpha}{t-\beta}=u;
\end{equation*}
the differential $\frac{dt}{\sqrt{R(t)}}$ is changed into a differential of the form $\frac{du}{\sqrt{R_{1}(u)}}$, where the polynomial $R_{1}(u)$ has only terms of even degree; because the relation between the values $u_{1}$, $u_{2}$ of $u$, which correspond to the values $t_{1}$, $t_{2}$ of $t$, is precisely $u_{1}+u_{2}=0$. Similarly, the rational function $\Fcal(t)$ is changed into a certain rational function $F(u)$, and the relation
\begin{equation*}
\Fcal(t_{1})+\Fcal(t_{2})=0
\end{equation*}
becomes
\begin{equation*}
F(u_{1})+F(u_{2})=0, \quad \text{that is} \quad F(u)+F(-u)=0.
\end{equation*}
It follows from this that $F(u)$ will be of the form $uF_{1}(u^{2})$, and it will suffice to set $u^{2}=x$ to reduce the differential
\begin{equation*}
\frac{u F_{1}(u^{2})du}{\sqrt{R_{1}(u)}}
\end{equation*}
to a differential of the form (1). We can therefore state the following theorem:

\begin{theorem}
Let $\left[t,\frac{Nt-M}{Lt-N}\right]$ be a substitution of period 2 permuting the four roots $a, b, c, d$ of the equation $R(t)=0$ in pairs, and let $\Fcal(t)$ be a rational function such that one identically has
\begin{equation*}
\Fcal(t)+\Fcal\left(\frac{Nt-M}{Lt-N}\right)=0;
\end{equation*}
the integral
\begin{equation*}
\int\Fcal(t)\frac{dt}{\sqrt{R(t)}}
\end{equation*}
is a pseudo-elliptic integral.
\end{theorem}
This result was obtained in another way by M. Raffy (see Bulletin, t. XII, p. 51); the above process furthermore gives the means of effecting the reduction. The preceding statement still applies when $R(t)$ is of the third degree, provided one regards one of the four roots $a, b, c, d$ as infinite.

\bigskip

\hrule

{
  \color{blue}
\medskip

\begin{definition}[Involution]
An \textbf{involution} is a function $S$ that is its own inverse; that is, $S(S(t)) = t$ for all $t$ in its domain. In the context of this theorem, we consider an involution that is a non-trivial \textit{linear fractional transformation}, having the form:
$$ S(t) = \frac{at+b}{ct+d}, \quad ad-bc \neq 0 $$
Such a transformation always has two fixed points (which may coincide or include $\infty$), given by the solutions to $S(t)=t$.
\end{definition}

\begin{theorem*}[thm. 1. Goursat, 1887]
Let $R(t)$ be a quartic polynomial with distinct roots, and let $S(t)$ be a non-trivial involution that permutes the roots of $R(t)$. If the rational function $\Fcal(t)$ satisfies the anti-invariance condition $\Fcal(t) + \Fcal(S(t)) = 0$, then the integral
$$ I = \int \frac{\Fcal(t)}{\sqrt{R(t)}} dt $$
is elementary.
\end{theorem*}

\begin{proof}
Let the two fixed points of the involution $S(t)$ be $\alpha$ and $\beta$. We introduce the change of variable
$$ u = \frac{t - \alpha}{t - \beta} $$
Under this transformation, the involution $S$ on $t$ corresponds to the simple reflection $u \mapsto -u$. Let us examine the effect of this substitution on the integrand.

The rational function $\Fcal(t)$ transforms into a new rational function $F(u) = \Fcal(t(u))$. The condition $\Fcal(t) + \Fcal(S(t)) = 0$ becomes $F(u) + F(-u) = 0$, which implies that $F(u)$ is an odd function of $u$.

The polynomial $R(t)$ transforms into a new expression $R_1(u)$ under the radical. Since $S$ permutes the roots of $R(t)$, the mapping $u \mapsto -u$ must permute the roots of $R_1(u)$. A polynomial whose roots occur in opposite pairs must be an even function. Thus, $R_1(u)$ is a polynomial in $u^2$, which we denote by $R_2(u^2)$.

The differential form $\frac{dt}{\sqrt{R(t)}}$ can likewise be shown to transform into an even function of $u$. The integral therefore becomes
$$ I = \int F(u) \cdot (\text{even function of } u) \, du = \int (\text{odd function of } u) \, du $$
The resulting integrand must have the structure $\frac{u \cdot H(u^2)}{\sqrt{R_2(u^2)}}$ for some rational function $H$.

We now apply a second substitution:
$$ x = u^2 \implies dx = 2u\,du $$
The integral is reduced to:
$$ I = \int \frac{H(x)}{2\sqrt{R_2(x)}} dx $$
Since $R(t)$ is a quartic, $R_1(u)$ is also a quartic, and therefore $R_2(x)$ must be a quadratic polynomial in $x$. The final integral is of the form $\int G(x) / \sqrt{ax^2+bx+c} \, dx$ for a rational function $G(x)$. Such an integral is known to be elementary.
\end{proof}
}
\medskip

\hrule

\bigskip

\section*{3}
The four roots $a, b, c, d$ being arbitrary, there exist three linear substitutions of period 2, which permute these roots in pairs. For example, the substitution
\begin{equation*}
\left\{t,\frac{(ab-cd)t+(a+b)cd-(c+d)ab}{[(a+b)-(c+d)]t-(ab-cd)}\right\}
\end{equation*}
permutes the two roots $a$ and $b$ and the two roots $c$ and $d$. The other two substitutions are deduced from this one, by successively taking any two of the roots for $a$ and $b$, and the other two for $c$ and $d$. Let us denote by $S_{1}$, $S_{2}$, $S_{3}$ these three substitutions; with the identical substitution $I(t)=t$, they form a group of four substitutions, and we have the relations
\begin{align*}
S_{1}^{2}=S_{2}^{2}=S_{3}^{2}&=I, \\
S_{1}S_{2}=S_{2}S_{1}=S_{3}, \\
S_{1}S_{3}=S_{3}S_{1}=S_{2}, \\
S_{2}S_{3}=S_{3}S_{2}&=S_{1}.
\end{align*}
Let $S\left[t,\frac{Nt-M}{Lt-N}\right]$ be in general a linear substitution and $F(t)$ a rational function of $t$; I will denote, for short, by $F(S)$ the rational function $F\left(\frac{Nt-M}{Lt-N}\right)$.

Given any polynomial of the fourth degree $R(t)$, we know from the above of three types of pseudo-elliptic integrals, relative to this polynomial,
\begin{equation*}
\int f(t)\frac{dt}{\sqrt{R(t)}}, \quad \int\varphi(t)\frac{dt}{\sqrt{R(t)}}, \quad \int\psi(t)\frac{dt}{\sqrt{R(t)}},
\end{equation*}
the rational functions $f(t)$, $\varphi(t)$, $\psi(t)$ respectively satisfying the relations
\begin{equation*}
f(S_{1})+f=0, \quad \varphi(S_{2})+\varphi=0, \quad \psi(S_{3})+\psi=0,
\end{equation*}
which are equivalent to these:
\begin{equation*}
f(S_{2})+f(S_{3})=0, \quad \varphi(S_{1})+\varphi(S_{3})=0, \quad \psi(S_{1})+\psi(S_{2})=0.
\end{equation*}
It is clear that by adding these three integrals, we still get a pseudo-elliptic integral
\begin{equation*}
\int\Fcal(t)\frac{dt}{\sqrt{R(t)}};
\end{equation*}
where $\Fcal(t)=f(t)+\varphi(t)+\psi(t)$. The preceding relationships, which are satisfied by the functions $f, \varphi, \psi$, give us, for the function $\Fcal(t)$, the relation
\begin{equation}
\Fcal+\Fcal(S_{1})+\Fcal(S_{2})+\Fcal(S_{3})=0.
\end{equation}
Conversely, let $\Fcal(t)$ be a rational function satisfying this condition (5). Let us set
\begin{equation*}
f=\frac{\Fcal-\Fcal(S_{1})+\Fcal(S_{2})-\Fcal(S_{3})}{4}, \quad \varphi=\frac{\Fcal+\Fcal(S_{1})-\Fcal(S_{2})-\Fcal(S_{3})}{4}, \quad \dots
\end{equation*}
We can therefore state the following theorem:

\begin{theorem}
Let $S_{1}$, $S_{2}$, $S_{3}$ be the three linear period 2 substitutions that permute the four roots of the equation $R(t)=0$ in pairs, and let $\Fcal(t)$ be a rational function such that one identically has
\begin{equation*}
\Fcal+\Fcal(S_{1})+\Fcal(S_{2})+\Fcal(S_{3})=0;
\end{equation*}
the integral
\begin{equation*}
\int\Fcal(t)\frac{dt}{\sqrt{R(t)}}
\end{equation*}
is a pseudo-elliptic integral.
\end{theorem}

\bigskip

\hrule

\medskip

{
\color{blue}

\begin{theorem*}[Goursat's Second Theorem]
Let $R(t)$ be a quartic polynomial with distinct roots, and let $S_1, S_2, S_3$ be the three distinct involutions that permute the roots of $R(t)$. If a rational function $\Fcal(t)$ satisfies the summation condition
$$ \Fcal(t) + \Fcal(S_1(t)) + \Fcal(S_2(t)) + \Fcal(S_3(t)) = 0 $$
then the integral $I = \int \frac{\Fcal(t)}{\sqrt{R(t)}} dt$ is elementary.
\end{theorem*}

\begin{proof}
The set of transformations $G = \{I, S_1, S_2, S_3\}$, where $I$ is the identity map, forms a group isomorphic to the Klein four-group ($V_4$). The group properties include $S_1^2=I$, $S_1 S_2 = S_3$, $S_2 S_1 = S_3$, etc.

The core of the proof is to decompose $\Fcal(t)$ into a sum of functions, each of which satisfies the simpler condition of Theorem 1. We can project $\Fcal(t)$ onto the subspaces corresponding to the non-trivial characters of $V_4$.

Let us define three components of $\Fcal(t)$:
\begin{align*}
\Fcal_1(t) &= \frac{1}{4} \left( \Fcal(t) - \Fcal(S_1 t) - \Fcal(S_2 t) + \Fcal(S_3 t) \right) \\
\Fcal_2(t) &= \frac{1}{4} \left( \Fcal(t) - \Fcal(S_1 t) + \Fcal(S_2 t) - \Fcal(S_3 t) \right) \\
\Fcal_3(t) &= \frac{1}{4} \left( \Fcal(t) + \Fcal(S_1 t) - \Fcal(S_2 t) - \Fcal(S_3 t) \right)
\end{align*}
The summation condition given in the theorem statement, $\Fcal + \Fcal(S_1) + \Fcal(S_2) + \Fcal(S_3) = 0$, ensures that $\Fcal(t)$ is the sum of these three components:
$$ \Fcal(t) = \Fcal_1(t) + \Fcal_2(t) + \Fcal_3(t) $$
Therefore, the integral can be split into three parts:
$$ I = \int \frac{\Fcal_1(t)}{\sqrt{R(t)}} dt + \int \frac{\Fcal_2(t)}{\sqrt{R(t)}} dt + \int \frac{\Fcal_3(t)}{\sqrt{R(t)}} dt $$
The proof is complete if we can show that each of these three integrals is elementary. We do this by showing that each component $\Fcal_i$ satisfies the anti-invariance condition of Theorem 1 for at least one involution.

Let's check $\Fcal_1(t)$ for anti-invariance under $S_1$:
\begin{align*}
4(\Fcal_1(t) + \Fcal_1(S_1 t)) &= \left( \Fcal - \Fcal S_1 - \Fcal S_2 + \Fcal S_3 \right) + \left( \Fcal S_1 - \Fcal S_1^2 - \Fcal S_2 S_1 + \Fcal S_3 S_1 \right) \\
&= \left( \Fcal - \Fcal S_1 - \Fcal S_2 + \Fcal S_3 \right) + \left( \Fcal S_1 - \Fcal - \Fcal S_3 + \Fcal S_2 \right) \\
&= 0
\end{align*}
Thus, $\Fcal_1(t) + \Fcal_1(S_1 t) = 0$. By Theorem 1, the integral $\int \frac{\Fcal_1(t)}{\sqrt{R(t)}} dt$ is elementary.

Similarly, we can show:
\begin{itemize}
    \item $\Fcal_2(t)$ is anti-invariant under $S_2$: $\Fcal_2(t) + \Fcal_2(S_2 t) = 0$.
    \item $\Fcal_3(t)$ is anti-invariant under $S_3$: $\Fcal_3(t) + \Fcal_3(S_3 t) = 0$.
\end{itemize}
Therefore, the integrals for $\Fcal_2$ and $\Fcal_3$ are also elementary by Theorem 1.

Since the original integral $I$ is a sum of three elementary integrals, it is itself elementary.
\end{proof}
}

\medskip
      
\hrule

\bigskip

The theorem still applies when $R(t)$ is of the third degree, provided we look at one of the roots $a, b, c, d$ as infinite. The sum $\Fcal+\Fcal(S_{1})+\Fcal(S_{2})+\Fcal(S_{3})$ is a symmetric function of the quantities $a, b, c, d$; one will be able to recognize if it is null without knowing the roots of the polynomial under the radical.

\subsection*{Examples}
Let $R(t)=(1-t^{2})(1-k^{2}t^{2})$. The integral
\begin{equation*}
\int f(t)\frac{dt}{\sqrt{(1-t^{2})(1-k^{2}t^{2})}}
\end{equation*}
will be pseudo-elliptic, provided that we have
\begin{equation*}
f(t)+f(-t)+f\left(\frac{1}{kt}\right)+f\left(-\frac{1}{kt}\right)=0.
\end{equation*}
The same will be true of the integral
\begin{equation*}
\int\frac{f(x^{2})dx}{\sqrt{(1-x^{2})(1-k^{2}x^{2})}},
\end{equation*}
provided that one has
\begin{equation*}
f(x^{2})+f\left[\frac{1-k^{2}x^{2}}{k^{2}(1-x^{2})}\right]+f\left(\frac{1}{k^{2}x^{2}}\right)+f\left(\frac{1-x^{2}}{1-k^{2}x^{2}}\right)=0.
\end{equation*}
(See HERMITE, Journal de Liouville, 1880; RAFFY, loc. cit., p. 71.)

\section*{4}

The preceding analysis shows us a relationship between some pseudo-elliptic integrals and the linear substitutions that permute the roots of the equation $R(t)=0$. If the quantities $a, b, c, d$ are arbitrary, there are no other linear substitutions than the three previous ones that permute these four roots. But, for some special forms of $R(t)$, there will be other substitutions of this kind. It is natural to wonder whether one can also attach pseudo-elliptic integrals to these new substitutions. This is indeed what happens.

Several assumptions can be made about these substitutions. Let $a, b, c, d$ always be the roots of $R(t)=0$; it may be that there is a linear substitution that does not change the point $a$ and permutes the other three roots $b, c, d$ circularly. This substitution will necessarily be of period 3.

Let $S$ be such a substitution. Now let $F(t)$ be a rational function of $t$ having the property expressed by the equation
\begin{equation*}
F(S)=\alpha F \quad \text{where} \quad \alpha=e^{\frac{2i\pi}{3}},
\end{equation*}
the integral
\begin{equation*}
\int F(t)\frac{dt}{\sqrt{R(t)}}
\end{equation*}
is pseudo-elliptic. Let's first do the transformation $z = \frac{t-a}{t-a_1}$, where $a, a_1$ are the fixed points of $S$. The transform of $S$ by this substitution is $S^{\prime}[z; \alpha z]$. The differential $\frac{dt}{\sqrt{R(t)}}$ turns into a differential of the form
\begin{equation*}
\frac{dz}{\sqrt{z(Az^{3}+B)}},
\end{equation*}
since the three non-zero roots of the new polynomial under the radical must be circularly permuted by the substitution $S'$. The rational function $F(t)$ changes into a rational function $F_1(z)$ verifying the relationship
\begin{equation*}
F_{1}(\alpha z)=\alpha F_{1}(z),
\end{equation*}
which will, therefore, be in the form
\begin{equation*}
F_{1}(z)=z\varphi(z^{3}).
\end{equation*}
The preceding differential becomes
\begin{equation*}
\frac{\varphi(z^{3})z\,dz}{\sqrt{z(Az^{3}+B)}}.
\end{equation*}
It now suffices to set $z^{3}=x$ to be reduced to a differential of the form (1); from which results the theorem stated above.

One can link to this case of reduction a rather famous integral, studied by Legendre and Clausen. If there is a linear substitution of period 3 leaving one of the roots $a, b, c, d$ fixed and permuting the others circularly, one can always, by a linear substitution, bring $R(t)$ to the form $R(t)=t^{3}-1$; which amounts to supposing $a=\infty$, $b=1$, $c=\alpha$, $d=\alpha^{2}$. There are eight period 3 substitutions, which permute three of these roots and leave the fourth fixed. To the four substitutions fixing one finite root correspond four types of pseudo-elliptic integrals. The substitutions to reduce these integrals are respectively
\begin{equation*}
t^{3}=x, \quad \left(\frac{t-1}{t+2}\right)^{3}=x, \quad \left(\frac{t-\alpha}{t+2\alpha}\right)^{3}=x, \quad \left(\frac{t-\alpha^{2}}{t+2\alpha^{2}}\right)^{3}=x.
\end{equation*}
The first is immediate. As for the second, an easy calculation allows one to verify that by setting $x=\left(\frac{t-1}{t+2}\right)^{3}$, we have
\begin{equation*}
\frac{1}{3}\frac{dx}{\sqrt{x(1-x)}}=\frac{t-1}{t+2}\frac{dt}{\sqrt{t^{3}-1}}.
\end{equation*}
Let A, B, C be three arbitrary constants; the integral
\begin{equation*}
\int\left(A\frac{t-1}{t+2}+B\frac{t-\alpha^{2}}{t+2\alpha^{2}}+C\frac{t-\alpha}{t+2\alpha}\right)\frac{dt}{\sqrt{t^{3}-1}}
\end{equation*}
will be expressed in finite terms using elementary symbols. This integral can be written
\begin{equation*}
\int\frac{\lambda(t^{3}-4)+\mu t^{2}+\nu t}{t^{3}+8}\frac{dt}{\sqrt{t^{3}-1}}
\end{equation*}
where $\lambda=A+B+C$, $\mu=-3(A+B\alpha^{2}+C\alpha)$, $\nu=6(A+B\alpha+C\alpha^{2})$. The coefficients $\lambda, \mu, \nu$ are, like A, B, C, arbitrary. Therefore the integrals
\begin{equation*}
\int\frac{(t^{3}-4)dt}{(t^{3}+8)\sqrt{t^{3}-1}}, \quad \int\frac{t^2 dt}{(t^{3}+8)\sqrt{t^{3}-1}}, \quad \int\frac{t\,dt}{(t^{3}+8)\sqrt{t^{3}-1}}
\end{equation*}
are pseudo-elliptic integrals; the last one is precisely the integral of Legendre and Clausen.

\section*{5}
The linear substitutions which permute the four quantities $\infty, 1, \alpha, \alpha^{2}$ form a group of twelve substitutions isomorphic to the group formed by the rotations which bring a regular tetrahedron onto itself.
Let $S_{i}$ be one of the twelve substitutions and $\mu_{i}$ the corresponding multiplier. Let us take a rational function $F(t)$ satisfying the two relations
\begin{equation}
\begin{cases}
\sum_{i=1}^{12}F(S_{i})=0, \\
\sum_{i=1}^{12}\frac{1}{\mu_i}F(S_{i})=0;
\end{cases}
\end{equation}
the integral
\begin{equation*}
\int F(t)\frac{dt}{\sqrt{t^{3}-1}}
\end{equation*}
is a pseudo-elliptic integral. To demonstrate this, I set
\begin{align*}
F(t)+F\left(\frac{t+2}{t-1}\right)+F\left(\frac{t+2\alpha}{\alpha^{2}t-1}\right)+F\left(\frac{t+2\alpha^{2}}{\alpha t-1}\right) &= 4\varphi(t), \\
F(t)-F\left(\frac{t+2}{t-1}\right)+F\left(\frac{t+2\alpha}{\alpha^{2}t-1}\right)-F\left(\frac{t+2\alpha^{2}}{\alpha t-1}\right) &= 4\varphi_{1}(t), \\
F(t)-F\left(\frac{t+2}{t-1}\right)-F\left(\frac{t+2\alpha}{\alpha^{2}t-1}\right)+F\left(\frac{t+2\alpha^{2}}{\alpha t-1}\right) &= 4\varphi_{2}(t), \\
F(t)+F\left(\frac{t+2}{t-1}\right)-F\left(\frac{t+2\alpha}{\alpha^{2}t-1}\right)-F\left(\frac{t+2\alpha^{2}}{\alpha t-1}\right) &= 4\varphi_{3}(t);
\end{align*}
from which we deduce
\begin{equation*}
F(t)=\varphi(t)+\varphi_{1}(t)+\varphi_{2}(t)+\varphi_{3}(t).
\end{equation*}
But the functions $\varphi_{1}, \varphi_{2}, \varphi_{3}$ respectively satisfy the relations
\begin{equation*}
\varphi_{1}\left(\frac{t+2}{t-1}\right)+\varphi_{1}(t)=0, \quad \varphi_{2}\left(\frac{t+2\alpha}{\alpha^{2}t-1}\right)+\varphi_{2}(t)=0, \quad \varphi_{3}\left(\frac{t+2\alpha^{2}}{\alpha t-1}\right)+\varphi_{3}(t)=0,
\end{equation*}
so that, according to the proposition of n°2, the three integrals
\begin{equation*}
\int\varphi_{1}(t)\frac{dt}{\sqrt{t^{3}-1}}, \quad \int\varphi_{2}(t)\frac{dt}{\sqrt{t^{3}-1}}, \quad \int\varphi_{3}(t)\frac{dt}{\sqrt{t^{3}-1}}
\end{equation*}
are pseudo-elliptic integrals. We still have to consider the integral $\int\varphi(t)\frac{dt}{\sqrt{t^{3}-1}}$, but, taking into account the relations (6), we will have
\begin{equation*}
\varphi(t)+\varphi(\alpha t)+\varphi(\alpha^{2}t)=0, \quad \varphi(t)+\alpha^{2}\varphi(\alpha t)+\alpha\varphi(\alpha^{2}t)=0,
\end{equation*}
so that $\varphi(t)$ will be of the form $\varphi(t)=t^{2}\psi(t^{3})$, and the considered differential takes the form
\begin{equation*}
t^{2}\psi(t^{3})\frac{dt}{\sqrt{t^{3}-1}},
\end{equation*}
which one reduces immediately to the form (1), by setting $t^{3}=x$.

\section*{6}

One can still make two hypotheses on the linear substitutions which permute the four roots $a, b, c, d$:
\begin{enumerate}
    \item It may be that a substitution, having for its double points two of these roots, permutes the other two.
    \item It may happen that there exists a linear substitution, circularly permuting the four roots.
\end{enumerate}
In the first case, the differential $\frac{dt}{\sqrt{R(t)}}$ can be reduced, by a linear substitution, to the form $\frac{dt}{\sqrt{t(t^{2}+1)}}$, and, in the second case, to the form $\frac{dt}{\sqrt{t^{4}-1}}$. Besides, one passes from the first to the second by changing $t$ into $\frac{1-t}{1+t}$. It is therefore sufficient to consider the second form. There exist eight substitutions permuting the four quantities $+1, -1, +i, -i$; they are the following:
\begin{gather*}
[t;t], \quad [t;it], \quad [t;-t], \quad [t;-it], \\
\left[t;\frac{1}{t}\right], \quad \left[t;\frac{i}{t}\right], \quad \left[t;-\frac{1}{t}\right], \quad \left[t;-\frac{i}{t}\right];
\end{gather*}
they form a group isomorphic to the group formed by the rotations which bring a system of two rectangular diameters of the sphere back upon itself. One can make correspond to these substitutions the following types of reducible integrals:
\begin{equation*}
\int t \varphi(t^2) \frac{dt}{\sqrt{t^4-1}}, \quad \int (t \mp \frac{1}{t}) \varphi(t \pm \frac{1}{t}) \frac{dt}{\sqrt{t^4-1}}, \quad \dots
\end{equation*}
but it should be noted that the substitutions of period 4 do not provide new pseudo-elliptic integrals.

\section*{7}

Given a polynomial $R(t)$ of degree higher than the fourth, having only simple linear factors, there does not in general exist a linear substitution permuting the roots of the equation $R(t)=0$. If such substitutions exist, one can make correspond to them ultra-elliptic integrals of the form
\begin{equation*}
\int f[t,\sqrt{R(t)}]dt,
\end{equation*}
which it is possible to reduce, by an algebraic substitution, to integrals of a lower genus. Without entering into all the details here, I will recall the two following examples, which are well known.

If one is dealing with a sixth-degree polynomial whose roots are in involution, one knows that one can, by a linear substitution, reduce it to the form
\begin{equation*}
R(u)=au^{6}+bu^{4}+cu^{2}+d,
\end{equation*}
and the two integrals of the first and second kind
\begin{equation*}
\int\frac{du}{\sqrt{R(u)}}, \quad \int\frac{u^2 du}{\sqrt{R(u)}}
\end{equation*}
are reduced by the substitution $u^{2}=x$ to elliptic integrals. This is basically the case of reduction found by Jacobi (Crelle's Journal, t. 8).

As another example, let us take $R(t)=t(t^{4}-1)$; the six quantities $0, \infty, \pm1, \pm i$ are permuted by a group of twenty-four linear substitutions, derived from the two substitutions
\begin{equation*}
[t;it], \quad \left[t;\frac{1-t}{1+t}\right];
\end{equation*}
this group is isomorphic to the group formed by the rotations which bring a regular octahedron back upon itself. The substitutions of period 2, which permute the roots in pairs, are four in number:
\begin{equation*}
\left[t;\frac{1-t}{1+t}\right], \quad \left[t;\frac{1+t}{1-t}\right], \quad \dots
\end{equation*}
To each of them correspond two integrals of the first kind, having only two distinct periods,
\begin{equation*}
\int\frac{(t\pm1\pm\sqrt{2})dt}{\sqrt{t(t^{4}-1)}}, \quad \dots
\end{equation*}
According to a theorem due to M. Poincaré, one concludes from this that there will be an infinity of integrals of the form
\begin{equation*}
\int\frac{(\alpha t+\beta)dt}{\sqrt{t(t^{4}-1)}}
\end{equation*}
having only two distinct periods.

\end{document}
